\chapter{What the Ordering Cannot Afford}\label{ch:cannot}

An ordering is a budget of a different kind: it says which rewrites are affordable, and a rule the
ordering cannot afford is not a rule the engine can have. This chapter is about three rules that
were wanted and could not be afforded as rewrites --- expansion, radical simplification, and one
trigonometric identity --- and about the three different things that were done instead. Each of
them has a lesson that the ordering chapter alone would not teach.

\section{Expansion is a function, not a rule set}

Distribution $a(b + c) \to ab + ac$ grows a term, and it grows it in a way no interpretation
compatible with $x \cdot x \to x^2$ can pay for. The proof that it cannot is short and was the
first thing tried: an interpretation under which distribution decreases makes the product $ab$
cheaper than the sum $a + b$ scaled by the multiplier, and that same interpretation makes
$x \cdot x \to x^2$ increase. It is the wall of Chapter~\ref{ch:order} again, seen from the other
side.

So expansion is not a rewrite system. \lc{Expand.dist} is one total, structurally recursive
function: children first, then a product with a sum among its factors is multiplied out, collecting
like monomials as the product is built factor by factor. The pipeline of Chapter~\ref{ch:order}
then collects whatever is left.

\begin{minted}{lean}
def distMul (fs : List Expr) : Expr :=
  if fs.any isAdd
    then ofPoly ((fs.map toPoly).foldl polyMul [⟨Q.one, []⟩])
    else .mul fs
\end{minted}

The collection during distribution is not a nicety. The first version distributed and left the
collection to the pipeline's collect-like-terms rule, which took the 4096 monomials of
$(x+y)^{12}$ one pair at a time at the root, with a full canonical sort after each pair: minutes,
and a recorded derivation of 49\,143 steps. Collecting as a polynomial --- sorted keys, coefficients
folded --- makes it instant and cuts the recorded expansion to two steps.

The soundness theorem is unconditional, because distribution holds in every commutative ring:

\begin{minted}{lean}
theorem dist_sound (ρ : EnvR) : ∀ e : Expr, evalR ρ (Expand.dist e) = evalR ρ e
\end{minted}

It is proved once for an abstract semantics --- a structure \lc{Sem} with an evaluator and its
four equations for sums, products, numerals and natural powers --- and instantiated twice, for
\lc{evalR} and for \lc{evalD}. The second instance matters in Chapter~\ref{ch:integrate}.

\begin{logbox}[Two shapes that read the same]
  A step-recording algorithm and a rewrite system look identical in the notebook: a list of
  steps, each with a rule name and a before and after. That is the argument for writing a
  function whenever a joint ordering would be forced. The reader loses nothing; the proof
  becomes an induction on the function instead of a fold over rules; and the fuel rewriter,
  which expansion had been the last user of, is deleted. There is no step budget anywhere in the
  engine except one, and that one is Chapter~\ref{ch:worlds}'s.
\end{logbox}

\section{Radicals: what the ordering allows}

Alex asked for $\sqrt 8 \to 2\sqrt 2$, and asked for it by retuning $M$. The analysis showed that
retuning cannot do it, and the log records both the request and the outcome.

As a term, $2\sqrt 2$ is $2 \cdot 2^{1/2}$, whose weight is $19$; $8^{1/2}$ weighs $11$. The result
contains the input's radical \emph{and more}. Only a weight that depends on the \emph{value} of a
numeral could pay for that --- making $8$ heavier than $2$ by enough --- and numeral weights must
stay bounded, or fold-constants stops decreasing on arbitrary sums ($1 + 7 \to 8$ would gain).

What the ordering does allow is the single-power form: $8^{1/2} \to 2^{3/2}$, at equal $M$ and
equal size. To order it, a sixth tier was added: the magnitudes of the integer numerals, counted
\emph{head-only}. That last word cost a day. A tier that reads a child's value fails the
children-monotonicity lemma the rewriter rests on --- a numeral child of equal weight could be
replaced by another --- and putting the value on the numeral node itself sidesteps it. The other
constraint is that $4^{2/3} \to 2^{4/3}$ ties every bit-length or value weight on the exponent
side, which is why the tier weighs integers only and non-integer exponents nothing.

Two more rules then do what Mathematica does where it matters. Radicals with the same square-free
part collect in a sum ($\sqrt{50} - \sqrt{18} \to 2\sqrt 2$; a sum of two radicals outweighs one
product, so this decreases $M$), and same-index radicals multiply under one root ($\sqrt{12} \cdot
\sqrt 3 \to \sqrt{36} \to 6$). Each rule guards itself with the decidable decrease its proof needs,
so the ordering lemma is a split on the guard; the guards never fail in practice and are honest
boundaries, not fuel. All three are proved sound over $\mathbb{R}$ unconditionally, their bases
being positive integers.

The printer then shows the single-power form the textbook way: $2^{3/2}$ as $2\sqrt 2$,
$12^{1/2}$ as $2\sqrt 3$, $5\sqrt{12}$ as $10\sqrt 3$. Text output stays faithful to the term. The
notebook shows what a reader expects; the engine holds what the ordering can carry; and the
distance between them is one function in the printer.

\section{Sine over cosine is tangent}

The third rule was found by the integration checker (Chapter~\ref{ch:integrate}) refusing a
correct answer: $\int \tan x$ was rejected because $\frac{d}{dx}(-\ln \cos x)$ normalized to
$\sin x \cdot (\cos x)^{-1}$ and the engine did not identify that with $\tan x$. The rule is a new
case of \rn{simp.function}, it \emph{does} decrease --- one function node replaces a product of two
--- and it had to be proved in all three layers: the additive measure of Chapter~\ref{ch:rewrite},
the ordering of Chapter~\ref{ch:order}, and soundness over $\mathbb{R}$, which is unconditional
because Mathlib's $\tan$ is $\sin/\cos$ everywhere, junk values included.

The lesson is about the cost of a rule. Three layers of proof for one identity is the price of
every rule in the engine, and it is why rules are added when a checker refuses something and not
before.

\begin{trybox}
\begin{minted}{text}
expand((x + y)^12)
sqrt(50) - sqrt(18)
sqrt(12) * sqrt(3)
\end{minted}
  The first has two expansion steps; everything after them is the pipeline collecting. The second and third run the radical rules; click a step and
  see the theorem name.
\end{trybox}

\section{Exercises}

\begin{exercisebox}[The wall, from the other side]
  Suppose an interpretation $I$ makes $a(b+c) \to ab + ac$ strictly decrease for all $a, b, c$.
  Show that $I(x \cdot x) < I(x^2)$ cannot also hold with $I$ of the form in
  Chapter~\ref{ch:order}.
\end{exercisebox}

\begin{exercisebox}[Why head-only]
  Give a rule that is a decrease under a value-of-children tier and a term on which the
  children-monotonicity lemma ($\mu(\text{child}') \le \mu(\text{child}) \Rightarrow
  \mu(\text{parent}') \le \mu(\text{parent})$) fails for that tier.
\end{exercisebox}
